\chapter{The Fact}
\label{chap:the-fact}

\section*{What You Measure}
\label{se:what-you-measure}
\markboth{WHAT YOU MEASURE}{WHAT YOU MEASURE}

The dashboard says the car is moving at forty-two miles per hour.

That is the number the driver sees, and in ordinary speech there is nothing
wrong with saying that the speedometer measures speed. The phrase is convenient,
and most of the time it is good enough. But it hides the most important fact
about measurement: what the instrument records is not the same thing as what
the instrument is said to measure.

A modern vehicle does not usually begin with a little needle deciding how fast
the car is going. It begins closer to the wheel. A toothed ring turns with the
hub, and a small magnetic or Hall-effect sensor watches the teeth pass. Each
passing tooth produces an electrical pulse. The instrument counts those pulses
over time. That count is the record.

The speed displayed on the dashboard comes later. To turn pulses into miles per
hour, the system must assume that the tire has the circumference expected by
the calibration, that the wheel is rolling rather than slipping, that the teeth
on the ring are evenly spaced, and that the stored conversion from pulse count
to displayed speed is the right one. The dashboard number is not false because
it depends on these assumptions. It is useful precisely because those
assumptions are usually good enough.

But when one of them fails, the difference between the record and the
phenomenon becomes visible. Put larger tires on the car without recalibrating
the system, and the display underreads. Spin the wheels on ice, and the display
may rise while the car barely moves. Break a tooth off the ring, and the count
is systematically wrong. In each case, the sensor may be doing exactly what it
was built to do. The failure is not necessarily in the measurement. It may be in
the model that connects the measurement to speed.

A radar gun approaches the same vehicle from an entirely different direction.
It does not count wheel pulses. It emits a microwave signal at a known frequency
and reads the frequency of the signal reflected back from the moving car. The
returning frequency shifts in proportion to the vehicle's motion along the
beam. From that shift, the radar gun displays a speed.

Again, the instrument has not touched speed itself. It has recorded a frequency
difference. To turn that difference into miles per hour, it uses a model of
wave reflection and relative motion. If the radar is aimed directly along the
vehicle's path, the displayed number may agree closely with the dashboard. If
the radar is held at an angle, the reading is lower, because only the component
of motion along the beam contributes to the shift. At forty-five degrees, the
gun can be off by about thirty percent. The radar may be working perfectly. The
geometry is not.

GPS gives a third route. A receiver in the car listens to timing signals from
satellites. It compares arrival times, uses information about the satellites'
orbits, corrects for delays through the atmosphere, and solves for position.
Speed can then be computed from changes in position over time, or from shifts
in the satellite signals themselves. Even the clocks are part of the story:
without corrections for the way satellite clocks run at different rates than
ground clocks, a consequence of both their speed and their altitude, GPS
positions would drift by roughly ten kilometers a day.

So three instruments may display the same familiar number: forty-two miles per
hour. The dashboard, the radar gun, and the GPS receiver can all agree. But the
agreement is not trivial. One record is a pulse count from a wheel sensor. One
is a frequency shift in a reflected signal. One is a timing solution involving
satellite clocks and atmospheric corrections. The number may be the same; the
records are not.

When the instruments agree, it is tempting to say that they have all measured
the same speed. A better way to say it is that their different records, under
their different models, can be reconciled to the same coarser description. The
agreement is an achievement. It is the result of calibration, comparison, and
work. Speed is not sitting inside any one instrument waiting to be copied onto a
display. Speed is the invariant that survives when multiple records are
reconciled.

This also explains why silence can be a valid record. Load the car onto a
flatbed truck and carry it across the state. The car has moved in the ordinary
sense. But if the wheels do not rotate, the wheel sensor records no pulses. The
dashboard speedometer, if it is watching only that sensor, correctly reports
zero. A GPS receiver in the same car may record the whole trip. The two records
do not contradict each other. They are written in different alphabets.

Ordinary language smooths over this gap. We say the speedometer measures speed.
We say the radar gun measures speed. We say GPS measures speed. The phrases are
useful because everyone knows roughly what they mean. But useful vagueness is
still vagueness. It hides the route from record to reading, from reading to
phenomenon, and from one instrument's ledger to another's.

A formal language cannot afford that kind of blur. It cannot honestly say that
an instrument measures speed without also specifying what the instrument can
actually report, what rule connects the report to a displayed value, and what
must be true for that rule to apply. The gap has to be written down. The
paperwork has to travel with the claim.

That is the reason for the compiler in this book. The compiler is not here
because computers are fashionable, and it is not here because physics needs a
new metaphor. It is here because a compiler can be forced to keep the paperwork.
It can be made to accept only those marks whose conditions have been supplied.
It cannot be told, "you know what I mean." It does not know. It checks.

This book asks what happens when we try to build physics, the discipline that
most confidently claims to measure the world, inside a language that will not
let us confuse what was recorded with what the record is about.

The hard part is that ordinary language hides the separation almost as soon as
it appears. We say speed, not pulse count. We say speed, not Doppler shift. We
say speed, not satellite timing solution. The single dashboard word gathers the
record, the rule, the calibration, and the interpretation into one comfortable
noun. That comfort is useful for driving. It is dangerous for foundations.

A formal language pulls the pieces apart again. It asks what was recorded, what
rule was used, what claim is being made, and what evidence travels with the
claim. It does not let the word ``speed'' absorb all of those obligations at
once. If the dashboard reading is to be carried forward, the path from wheel
pulse to displayed value must be made explicit enough that the next step can
depend on it.

This is where the compiler becomes more than a tool for checking syntax. In
this book, the compiler is the instrument that watches the argument. It will
not move to the next mark because the previous sentence sounded plausible. It
will not forgive a missing condition because the metaphor was good. It accepts
only what has been built in the language it understands.

That does not make the compiler wise. It does not understand cars, radar, GPS,
or physics. Its strength is narrower and more useful: it can be made to refuse
a distinction that has not been supplied. It can keep a ledger of what has
actually entered the formal record. If a later construction needs an earlier
piece of evidence, the compiler must be able to find it.

Chapter 1 begins with the smallest such entry. Before speed, before clocks,
before particles, before geometry, and even before counting in the ordinary
sense, the device needs one accepted distinction. It needs a claim small enough
for the compiler to check and sturdy enough for later marks to stand on. It
begins with a fact.

% EPISODE COVERAGE: Episode 1 (partial) -- Fact, Number, CarrierProcess, le duality
% TONE: Conversational, wonder-forward. The compiler as a character with one obsession.

\markboth{THE FACT}{THE FACT}
\section{The Compiler as Character}
\label{se:compiler-as-character}

The first character in this book is not a physicist, not a particle, and not
the universe. It is a compiler.

That sounds like a cold place to begin, but a compiler is one of the most
careful instruments human beings have ever built. It has almost no imagination.
It does not admire a beautiful argument. It does not forgive an elegant mistake.
It does not care whether a sentence sounds plausible, famous, revolutionary, or
comforting. It asks a smaller question: does this term fit here?

Lean is a proof assistant. Its surface resembles a programming language, but
its central habit is closer to inspection than execution. Give Lean a
definition, and it checks whether the definition has the type claimed for it.
Give Lean an inference, and it checks whether the inference follows from the
context already accepted. Give Lean a structure with missing fields, and it
stops. Give Lean a theorem without a proof, and it stops. Give Lean a term whose
implicit assumptions cannot be supplied, and it stops.

This stopping is not failure in the ordinary sense. It is the measurement.

An ordinary compiler accepts a program when the syntax and types line up well
enough to produce executable code. Lean asks for more. It elaborates.
Elaboration is Lean's process of filling in all the details the author left
implicit: hidden arguments, typeclass instances, and the evidence required to
close each term. It unfolds definitions when permitted, checks that every
expression has the type it claims to have, and refuses to proceed when a
required piece of evidence is missing. The process is local and incremental.
Lean does not get to use tomorrow's
definition to justify today's declaration. It cannot appeal to what the author
intended to write three pages later. It carries only what has already been
accepted.

This is why Lean will serve as the instrument in the book. It behaves like a
ledger with a gate attached. Each new entry must be consistent with the entries
already recorded. Once accepted, an entry becomes part of the context in which
the next entry will be judged. The compiler is not asked to survey the whole
universe from above. It is asked to stand at the boundary of the record and
decide whether one more mark may be added.

That is also how measurement works.

A measuring instrument does not receive the world in full. A voltmeter does not
record the complete electromagnetic state of a room. A thermometer does not
store the positions and momenta of all the molecules whose motion it reports. A
camera does not capture every photon that could have arrived. Each instrument
has a grammar of admissible responses. It can say this, not that. It can record
one mark from a finite or countable set of marks. It can distinguish some
differences and remain silent about others.

The reading on the instrument is therefore not the world. It is the world's
passage through a rule.

Lean is like that. It does not know what the comments mean. The jokes, the
epigraphs, the arguments in the margins, the sudden claims about physics, the
little acts of theatrical bravado in the source file -- Lean ignores all of
them. The compiler cannot be charmed into accepting a bad term. It cannot be
threatened into skipping a missing field. The prose may strut, wink, or shout;
the elaborator reads the declarations.

This is important enough to say plainly: the comments are voice, but the terms
are the experiment.

The device we will follow was written as a sequence of Lean files called
episodes. The episodes are funny, unruly, and full of human weather. They talk
about John Henry, clocks, jokes, compilers, physics, and the danger of trusting
symbols too quickly. But underneath that voice, something very strict is
happening. The compiler is being asked to accept a tower of definitions, one
floor at a time. Each floor supplies a shape that the next floor will require.
Each accepted declaration becomes part of the environment. Each successful
instance says, in effect, that the promised structure can actually be inhabited.

The book you are reading is not a transcript of all that code. It is a guided
tour of why the acceptance of that code matters. The technical companion volume
develops the mathematical framework in full. Here, our concern is more direct:
what did the compiler have to agree to, step by step, before the final question
could even be asked?

The answer begins smaller than any physical law.

It begins before space, before time, before motion, before energy, before
fields, before particles, before geometry, and before arithmetic in the usual
sense. It begins with the compiler's most primitive discipline: a claim cannot
enter the record merely by being spoken. It must arrive in a form the compiler
can check.

The compiler's yes is not a compliment. It is a recorded compatibility.

This is the point at which the compiler becomes a character. Not because it has
desires, and not because it understands the story being told around it. It
becomes a character because it has a consistent temperament. It is strict,
literal, local, and patient. It never leaps ahead. It never accepts a promise in
place of a construction. It never reads the author's mind. It only asks whether
the present declaration can be closed against the current context.

There is something almost comic about that severity. The source file may make a
wild claim in a comment; Lean does not blink, because comments are not part of
the formal record. The next line may introduce a tiny definition; Lean examines
it with complete seriousness. The human author is free to talk about the
universe. The compiler asks whether the field has the right type.

This mismatch is part of the charm of the device, but it is also part of its
argument. If the compiler accepted the tower because it understood the jokes,
the experiment would be worthless. If it accepted the tower because it shared
the author's metaphors, the result would tell us nothing. The force of the
construction comes from the opposite fact: Lean does not know the story. It
knows only the terms.

The device therefore has two narratives running at once. One narrative is the
human one, full of analogy, provocation, and delight. The other is the compiler's
narrative, which is not a narrative at all but an ordered sequence of accepted
declarations. This book lives between them. It asks what a reader can learn
from the human story once the machine has already performed its colder check.

To see why this matters, imagine an instrument on a lab bench. A pulse enters.
The instrument either registers a mark or it does not. If it registers, the mark
is not merely decoration; it changes the state of the experiment. Future
readings must be interpreted in light of the mark already recorded. The
instrument has accumulated a history.

Lean accumulates context in the same way. A definition accepted earlier becomes
available later. A class introduced earlier can become a requirement for a
later structure. An instance supplied earlier can be found by the compiler when
a later declaration asks for it. The source file is not a bag of independent
sentences. It is an ordered record. The order matters because the compiler's
memory matters.

This is the first appearance of the ledger.

A ledger is not just a list. A shopping list can be rewritten. A ledger carries
the burden of what has already been entered. Once a bank records a deposit, the
next balance is computed with that deposit included. Once a laboratory notebook
records a reading, the next analysis cannot honestly pretend the reading was
never taken. Once Lean accepts a declaration, the next declaration is checked in
the world that acceptance has made.

The word ``world'' is still too large for what we have built. At the beginning
of Chapter 1, there is no world in the device. There is only the compiler's
environment and the possibility of adding something to it. But that is already
enough to pose the book's first question. If compilation is an ordered act of
acceptance, and measurement is an ordered act of recording, what is the smallest
thing that can be accepted as a measurement?

Not the smallest physical object. Not the smallest particle. Not the smallest
length. Those are much later questions, and they already assume an instrument
capable of reporting them. The first question is more primitive:

\begin{center}
\emph{What is the smallest thing the compiler can verify?}
\end{center}

The answer will look almost disappointingly simple. That disappointment is a
good sign. The first step should be boring. If the first step already contained
motion, matter, or geometry, the whole construction would be suspect. A book
that wants to explain measurement cannot begin by assuming the machinery of
measurement. It must begin with something smaller: a claim, and the evidence
that the claim can be decided.

In the next section, the device gives that smallest object a name.

\section{The Primitive Distinction}
\label{se:primitive-distinction}

The smallest object in the device is called a \lean{Fact}.

That name is almost too familiar. In ordinary speech, a fact is something true
about the world. The sun is hot. The cup is on the table. The detector clicked.
The word feels solid. It sounds like the end of an argument.

Lean needs something more careful than that.

The device does not begin by asking the compiler to hold the whole world. It
does not begin with a particle, a field, a clock, a ruler, or a law. It does
not even begin with a number. It begins with a claim that has been made precise
enough for the compiler to carry. That distinction matters. A claim by itself is
cheap. Anyone can make one. The device needs a claim with a way of settling
whether the claim holds.

Before the first code block, here is the little key needed for this chapter.
This is not a Lean lesson. It is only enough vocabulary to keep the first few
objects from looking like magic.

\begin{sidebar}{How to read the first Lean block}
\label{sb:how-to-read-lean}
\lean{class} names a promise. Anything called a \lean{Fact} must supply the
fields listed after \lean{where}.\\
\lean{Prop} is a proposition: a claim Lean can reason about.\\
\lean{Decidable p} is the receipt that the particular proposition \lean{p} can
be answered yes or no.\\
\lean{namespace} ... \lean{end} puts a definition under a shared name.
\end{sidebar}

The first promise is short.

\begin{leancode}
@[reducible]\\
class Fact where\\
\quad truth : Prop\\
\quad decTruth : Decidable truth
\end{leancode}

There are only two lines inside the promise. The first is
\lean{truth : Prop}. A \lean{Prop} is a proposition: a claim Lean can reason
about. It is not a simple switch, already flipped to yes or no. It is not the
answer to the question. It is the question stated in a form Lean can inspect.

This is worth slowing down for. The word \lean{truth} may tempt the reader to
think the field already contains truth in the ordinary sense. It does not. It
contains a claim. The claim might be as plain as "true is true." It might be
more complicated later. At this point, the compiler has only been handed
something that can be asked.

The second line is \lean{decTruth : Decidable truth}. This is the receipt. It
says that the claim in \lean{truth} does not arrive alone. It arrives with a
method for arriving at a definite answer. Not that the claim is true. Not that
every possible claim can be decided. Only this: for the proposition stored in
\lean{truth}, the device also carries evidence that the proposition can be
answered yes or no.

A \lean{Bool} would already be a value, already one side of the distinction. A
\lean{Prop} is only the claim. \lean{Decidable truth} is what lets this
particular claim cross into the ledger: the compiler is told that the claim can
be settled, not merely stated.

A smoke detector has this shape.

It does not know the whole room. It does not report the chemistry of every
particle in the air, the temperature of the ceiling, the history of the wiring,
or whether anyone remembered to make toast. It is built around a narrower
question: has the sensed condition crossed the alarm threshold? The world may
be complicated, but the detector is not trying to receive the world in full. It
is trying to make one distinction reportable.

The detector's threshold is its method for arriving at a definite answer. Above
the threshold, alarm. Below the threshold, no alarm. That answer may be crude.
It may depend on calibration, placement, battery condition, dust, humidity, and
the particular material burning. But the device is designed so that, under its
own rule, the question can be answered. The alarm is not a theory of fire. It is
a disciplined yes or no.

That yes or no is the beginning of a ledger. If the alarm sounds, the mark is
not merely an opinion. It is an accepted report made by an instrument according
to its construction. If the alarm does not sound while the detector is present
and operating, that silence is also meaningful: the threshold was not crossed
in a way the instrument could report. The detector's reading is therefore not
the whole truth of the room. It is a decided distinction.

A scale has a related shape. It does not report the complete mechanical history
of the object placed on it. It settles into a reading its gauge can display. The
reading depends on the instrument's construction, its calibration, and the
range of distinctions it can resolve. The scale does not tell us everything. It
tells us what its mechanism can make definite.

A wheel-speed sensor has the same shape, and this will become the book's
running example. The sensor does not record the car's position, fuel level,
road surface, or destination. Its rule is much smaller: a tooth on a rotating
ring has passed the sensor, or it has not. That pulse is not yet speed. It is a
single reportable mark from which speed may later be inferred.

So a \lean{Fact} is not merely a proposition. It is a proposition with a
receipt.

The proposition is the claim. The receipt is the attached evidence that the
claim can be settled. A loose claim may be interesting, inspiring, or useful,
but it is not yet a \lean{Fact} in this device. The compiler is being asked to
carry only claims that arrive with their paperwork.

The small decoration \lean{@[reducible]} above the promise tells Lean to look
through \lean{Fact} rather than treat it as a sealed box. The chapter does not
need the machinery behind that instruction. The reader only needs the posture:
the first object is meant to be transparent to the compiler. The device is not
hiding its opening move.

Once the promise exists, the device supplies the dullest possible example of
something that keeps it.

\begin{leancode}
namespace Fact\\
def Truth : Fact := \{\\
\hspace*{2em}truth := true,\\
\hspace*{2em}decTruth := Decidable.isTrue rfl\\
\}\\
end Fact
\end{leancode}

This defines \lean{Fact.Truth}: the dullest possible inhabitant of the first
promise. The claim is \lean{true}. The receipt is \lean{Decidable.isTrue rfl}.
\lean{rfl} is Lean's smallest shrug: of course true is true.

That is the point. The first accepted mark should not already contain physics,
geometry, or arithmetic. It should be small enough to inspect. \lean{Fact.Truth}
is a settled value, not a question waiting for an answer, and the compiler can
file it away for later use.

Notice the order. The device does not first build a number line, a clock, or a
space. It begins with the accepted mark itself. Only later will there be
counting. Only later will there be an instrument. Only later will the tower
acquire enough structure to talk about science, physics, and closure.

It may feel strange that so much attention is being paid to something so small.
But the smallness is the reason the object matters. We are trying to explain
why a large Lean device compiles successfully. The successful compile at the
end will not be a single miracle. It will be the last mark in a long series of
accepted marks. The first accepted mark therefore deserves care. If the first
step is vague, every later step inherits the blur.

The device refuses that blur. It says: begin with a claim, and carry the means
of deciding it.

This is also where the compiler begins to resemble an instrument. An instrument
does not merely announce reality. It is built to decide among admissible
readings. A voltmeter may show a voltage, but the displayed number is the end
of an enormous narrowing. The world pressed against the device in more ways
than the display can say. The gauge reported only what its construction made
reportable.

A \lean{Fact} is reportable in this sense. It has the shape of something that
can enter a ledger. The claim is there. The means of decision is there. The
compiler can carry both together.

That last phrase, "carry both together," is more important than it may first
appear. A claim without its receipt can drift away from the conditions under
which it was accepted. In ordinary speech, this happens constantly. Someone
remembers the conclusion but forgets the measurement. A number survives while
the calibration that produced it vanishes. A graph circulates without the
instrumental limits that made the data meaningful. The result may still look
like knowledge, but its support has been quietly lost.

The device does not allow that separation at the base. The proposition and the
evidence of decidability are packaged together. The fact is not merely the
mark. It is the mark with the authority to be a mark.

Here is the physical bridge. A measurement is not the whole world. It is an
accepted distinction. The thermometer distinguishes this range from that one.
The smoke detector distinguishes alarm from no alarm. The wheel-speed sensor
distinguishes a pulse from no pulse. Each does so according to a construction
that determines what may count as a report.

A \lean{Fact} is what such a report looks like when it has been made precise
enough for the compiler to carry it.

This does not mean the compiler has learned physics. Not yet. It means the
device has found a formal shape that resembles the smallest useful act of
measurement: a definite claim, together with the means by which the claim may
be decided. Later, the book will build instruments, ledgers, signals, clocks,
and transitions from this starting point. For now, the first lesson is narrower
and cleaner.

No receipt, no fact.

There is one more reason to keep this opening fact boring. Because the computer
running Lean is physical, even the accepted mark is embodied. The mark
recording this fact is held in place by electrons in silicon. Chapter 1 will
not chase that thread very far. It is enough to notice the irony: the first
object is abstract enough to fit in a proof assistant and physical enough to
live in a machine.

Counting requires this. We will see how shortly.


\section{Counting from Nothing}
\label{se:counting-from-nothing}

Once the device has a fact, it can begin to count. But the title of this
section is a trap. The device does not really count from nothing. Nothing, in
the strict sense, leaves no receipt. The first countable thing is not empty
space, silence in general, or an unmarked void. The first countable thing is a
fact that the compiler has accepted.

Here is the whole beginning of number in the episode.

\begin{leancode}
inductive Number where\\
\quad | zero : Fact -> Number\\
\quad | one  : Fact -> Number -> Number
\end{leancode}

The word \lean{inductive} means that the device is listing the possible ways to
build a value of this kind. It is a menu, but a closed menu. A \lean{Number} can
be built with \lean{zero}, or it can be built with \lean{one}. There is no
third constructor waiting in the margins. If a later line wants a
\lean{Number}, it must be able to get one by using these forms.

That is already a strong discipline. In ordinary thought, zero often feels
like absence. No apples. No pulses. No cars. But the device refuses to treat
absence as self-explanatory. The constructor

\begin{leancode}
zero : Fact -> Number
\end{leancode}

does not say that zero is an empty integer floating by itself. It says that the
base case of number is made from a \lean{Fact}. Even zero has paperwork.
Because \lean{Number} is inductive, Lean knows that every \lean{Number} must
eventually bottom out in a \lean{zero}: there is no way to build a
\lean{Number} that does not trace back to a base case.

Return to the wheel-speed sensor. Suppose the car is switched on, the sensor is
powered, and the toothed ring is being watched, but no tooth has yet passed the
pickup. The pulse count is zero. That zero is not the same as no instrument, no
wire, no rule, and no observation. It is a watched zero. It means that the
ledger has a place to stand before the first pulse arrives. The base case is
not silence; it is a standing receipt.

The second constructor is the act of adding one more accepted mark.

\begin{leancode}
one : Fact -> Number -> Number
\end{leancode}

Read it slowly. To make one more number, the device needs a fresh
\lean{Fact} and an earlier \lean{Number}. The new mark does not erase the old
count. It sits on top of it. The count grows by carrying its previous history
forward.

For the wheel-speed sensor, this is the passing tooth. One more tooth crosses
the pickup. The electronics register a pulse. But the pulse is not yet speed.
It is one accepted pulse, one more entry in a ledger that may later be
interpreted as speed under a model of wheel circumference, elapsed time, and
calibration. The count is still upstream of the displayed quantity.

A small number in this device therefore looks like a stack of receipts:

\begin{leancode}
one f3 (one f2 (zero f1))
\end{leancode}

The notation reads from the outside inward, but the history reads from the
inside outward. First there is \lean{zero f1}, the watched base case. Then
\lean{one f2} records one accepted successor over that base. Then
\lean{one f3} records another accepted successor over the result. The object
has three layers of formal evidence. It is not merely the numeral two. It is a
little ledger.

That is why the device does not begin with Lean's ordinary natural numbers.
Lean already has \lean{Nat}, and \lean{Nat} is enormously useful. But using it
here would hide the question the chapter wants to expose. The device is not
trying to borrow arithmetic. It is trying to build the precondition for
arithmetic: the possibility of repeated accepted marks.

So there is no addition yet. No subtraction. No multiplication. No comparison.
No number line in the ordinary sense. There is only a base entry and
a way to place one more witnessed entry over a previous one. That is enough to
begin counting, but not enough to do mathematics. The distinction matters. A
ledger that can make marks is older than the algebra that later teaches those
marks to move.

\begin{gphenom}{Peano-Kushim}
\label{ph:peano-kushim}
\GPhObs The oldest verified act of counting is a 5,000-year-old Sumerian tablet
recording a barley inventory, attributed to a scribe named Kushim.
Peano axiomatized the same act in 1889.
\GPhCon Counting requires a method for reaching a definite answer:
this item is distinct from that one.
The ledger begins with the first verified distinction, not with the first number.
\GPhSeq The compiler's \lean{Fact} is the minimum unit of counting:
a claim that can be verified. In this device, arithmetic is assembled from
verified distinctions, not assumed as given. One verified distinction at a time.
\end{gphenom}

The box matters because \lean{Number} is not a numeral first. It is an
entry-making device. Each step of the count carries a \lean{Fact}.

The compiler has now accepted that counting is possible. But a ledger that can
count marks without any rule for which marks are permissible is only half a
ledger. It can remember that another mark has arrived; it cannot yet say what
counts as a permissible continuation.

The next chapter will introduce a typeclass called \lean{ADMISSIBLE} that
formalizes exactly this constraint. The Hume box names it here as a warning:
the ledger cannot prophesy, only check.

\begin{gphenom}{Hume}
\label{ph:hume}
\GPhObs No finite sequence of observed sunrises guarantees that the sun will
rise tomorrow. Hume established this in 1739: induction cannot be
logically justified from finite evidence.
\GPhCon The compiler cannot infer, from $N$ verified facts, that the $(N+1)$th
fact will hold. Every general claim in the ledger is provisional.
\GPhSeq The ledger records what has been witnessed, not what must be.
The \lean{ADMISSIBLE} typeclass does not predict; it only checks
whether a proposed continuation is consistent with what is already recorded.
\end{gphenom}

Hume is here to mark a limit: a record of previous marks is not yet a license
for the next one.

Once the device can count, it can ask whether one count is smaller than
another. But the answer depends on something the counts carry with them.
Section 1.4 shows what.

\section{The Two Orderings}
\label{se:two-orderings}

Counting is still not comparison.

A wheel-speed sensor can count pulses as teeth pass the pickup. One pulse,
then another, then another. That gives the ledger a direction: the count grows
as marks are accepted. But the device has not built ordinary natural numbers.
Its counts are not bare lengths. Each layer carries a \lean{Fact}, and that
means comparison cannot look only at how many layers a number has. It must also
look at what those layers carry.

We could compare two numbers by counting layers, but that would ignore what the
layers carry. Two counts of the same depth may point in opposed directions. The
episode therefore defines its own ordering for \lean{Number}.

\begin{leancode}
namespace Number\\
def le : Number -> Number -> Prop\\
\quad | .zero \_ , \_ => True\\
\quad | .one \_ \_, .zero \_ => False\\
\quad | .one p1 n1', .one p2 n2' =>\\
\quad\quad match p1.decTruth, p2.decTruth with\\
\quad\quad | isTrue \_,  isTrue \_  => le n1' n2'\\
\quad\quad | isTrue \_,  isFalse \_ => False\\
\quad\quad | isFalse \_, isTrue \_  => True\\
\quad\quad | isFalse \_, isFalse \_ => \(\neg\) le n1' n2'\\
end Number
\end{leancode}

The type of this comparison is \lean{Number -> Number -> Prop}. That last word
matters. The comparison produces a proposition, a claim the compiler can reason
about, not a simple \lean{Bool} switch that automatically returns yes or no.
The device is still working in the logic layer. It is defining what it would
mean for one of these receipt-carrying numbers to be less than or equal to
another.

The first two cases are familiar. A \lean{zero} is less than or equal to
anything. A number built with \lean{one} is not less than or equal to
\lean{zero}. So far, this looks like the ordinary beginning of counting.

Then the recursive case changes the weather. When both numbers are built with
\lean{one}, the comparison opens the facts carried by the two outer layers.
The tails are not compared immediately. First the device asks how those two
facts decide.

\begin{center}
\begin{tabular}{lll}
left fact & right fact & comparison behavior \\
\hline
\lean{isTrue} & \lean{isTrue} & compare the tails in the same direction \\
\lean{isTrue} & \lean{isFalse} & reject the comparison \\
\lean{isFalse} & \lean{isTrue} & accept the comparison \\
\lean{isFalse} & \lean{isFalse} & negate the tail comparison
\end{tabular}
\end{center}

The first line is the direction one expects from ordinary counting. If both
outer facts are affirmative, the comparison peels one layer from each side and
continues with the tails. This is the direction of accumulation: one more mark,
then one more mark after that.

The two middle lines are fixed answers. If the left layer is affirmative and
the right layer is complementary, the comparison fails. If the left layer is
complementary and the right layer is affirmative, the comparison succeeds. The
device is treating the facts as sign-like or direction-like, but not as
ordinary integer signs stored in a separate field. The direction is carried by
the facts themselves.

The last line is the one to handle carefully. When both outer facts are
complementary, the code does not swap the two tails and call the comparison
again. It negates the proposition \lean{le n1' n2'}. In other words, the
statement that held in the forward direction no longer holds in the same way.
The negation is not a swap. The two tails are not exchanged. What held in one
direction is declared not to hold. That is a logical reversal, not a positional
one.

A countdown has this flavor. If a timer moves from ten seconds remaining to
five seconds remaining, one description says the remaining number has become
smaller. Another says the process has gone further. The same ticks can support
two readings, depending on which direction the ledger is asked to privilege.
A clock and its complement do something similar. One marks the beat where the
other marks the gap. Together they describe one cadence from opposite sides.

That is why the section is called the two orderings. The device has not proved
a grand theorem about covariance. It has defined a comparison relation whose
branches already contain two habits. One follows the buildup of accepted marks.
The other follows the way a complementary mark turns the reading around. In the
language the book will use later, these are the first shadows of covariant and
contravariant behavior: accumulation and unwinding.

The episode then registers this comparison under Lean's familiar ordering
symbol.

\begin{leancode}
instance : LE Number where\\
\quad le := Number.le
\end{leancode}

After this line, the compiler can read \lean{n1 <= n2} as shorthand for the
four-case rule above. The symbol looks familiar, but the meaning is local. It
belongs to this \lean{Number}, whose layers carry facts.

The compiler has not discovered physics. It has accepted a stranger and more
modest fact: once a count carries facts inside it, comparison is not merely a
matter of length. Direction has entered the ledger.

A count with direction is still not an instrument. Comparison can say how two
counts relate, but an instrument does not merely compare. It persists. It holds
a value between events, and it applies a rule when the next mark arrives.

\section{What the Compiler Carries Forward}
\label{se:what-compiler-carries}

A speedometer needs more than a pulse count. It needs something that carries
the mark, something that holds the current value, and some rule for what happens
when the next event is processed. A pulse without an update rule is only a
pulse. A count without a carrier is only a static ledger.

The final object in this chapter has that shape.

\begin{leancode}
structure CarrierProcess\\
\quad (Carrier: Type)\\
\quad where\\
\quad symbol: Fact\\
\quad value: Number\\
\quad event: Number -> Number := fun s =>\\
\quad\quad match s with\\
\quad\quad | .zero \_ => .zero Fact.Truth\\
\quad\quad | .one p \_ => .one p value
\end{leancode}

This is a \lean{structure}, a named package of named pieces. The name is
\lean{CarrierProcess}. The word \lean{Carrier} appears before the fields
because the process is abstract over the kind of thing doing the carrying. The
device has not yet chosen wheel teeth, radar waves, radio pulses, GPS timing
marks, or anything else. It is asking for the form of a process before it
chooses a physical medium.

The first field is:

\begin{leancode}
symbol: Fact
\end{leancode}

This is the carried mark. It is not optional. A \lean{CarrierProcess} cannot
exist without a \lean{symbol}, and that symbol has the shape Chapter 1 has been
building toward: a \lean{Fact}. The symbol may be dull. It may carry
\lean{Fact.Truth}. But it is present.

The second field is:

\begin{leancode}
value: Number
\end{leancode}

This is the current count or state held by the process. It is not an abstract
number line borrowed from elsewhere. It is one of the receipt-carrying numbers
from this chapter: a \lean{zero} with a fact, or a stack of \lean{one} layers,
each with its own fact.

The third field is:

\begin{leancode}
event: Number -> Number
\end{leancode}

This is the rule for transforming one number into another. With it, the process
is no longer merely a mark beside a count. It has a next-step rule. Given a
number, it can produce a number.

The source also supplies a default event. That default matters because it shows
how little machinery is needed before something memory-like appears. If the
incoming state is \lean{zero}, the rule returns \lean{zero Fact.Truth}, the
plainest accepted base case. If the incoming state is built with \lean{one}, the
rule keeps the outer fact \lean{p} and attaches it to the process's stored
\lean{value}. Suppose the process currently holds
\lean{value = one f2 (zero f1)}. If the incoming state is
\lean{one p (zero f3)}, the default event returns \lean{one p value}, that is,
\lean{one p (one f2 (zero f1))}. The incoming tail \lean{zero f3} is discarded.
The process does not add the incoming number. It stamps the incoming fact onto
its own stored history.

That is the fence-post edge of the construction. The next number is tied to the
value already being carried. In the speedometer image, the sensor does not lose
its count between wheel teeth. Between pulses, it holds its state. When another
event is processed, the update rule is interpreted against the value the
instrument already has.

This is still not a full physical instrument. Lean has not been asked to know
what a wheel is, what a road is, or what speed means. But the outline is now
recognizable. Something carries a witnessed mark. Something holds a count.
Something says how a count moves to a next count. The compiler has accepted the
shape of a minimal process.

That raises the question the Marconi box is meant to hold in place. If a process
can carry a mark and update a value, what happens when the instrument is
present and no event arrives? The answer cannot be that the symbol disappears.
The \lean{symbol} field is always present. Silence, if it matters, must be
recorded as a state of the watched instrument, not as the absence of the ledger
itself.

\begin{gphenom}{Marconi}
\label{ph:marconi}
\GPhObs Guglielmo Marconi's first transatlantic radio transmission in 1901
was three dots -- the letter S in Morse code. The receiver knew it was
a signal because it was distinguishable from silence. Silence itself
was informative: the channel was open.
\GPhCon Silence is an admissible ledger entry. The absence of a signal is
a witnessed distinction: the instrument was present and nothing arrived.
\GPhSeq A carrier must be able to represent both signal and silence as
first-class states. The symbol is always present; what varies is what
its \lean{Fact} records. A ledger that has no formal entry for
``nothing arrived'' cannot distinguish a watched channel from an
unobserved one.
\end{gphenom}

Marconi's lesson is not radio nostalgia. It is the difference between an
unwatched absence and a watched silence.

Chapter 1 has now assembled its first inventory. The compiler has accepted a
\lean{Fact}. It has accepted a number: a chain of receipt-carrying marks. It
has accepted an ordering: a rule that lets those chains be compared in more
than one direction. And it has accepted a carrier process: a mark, a value, and
a rule for what comes next.

The compiler has still not been asked to believe in cars, clocks, particles, or
fields. But it has accepted enough structure to carry a process forward. The
next chapter can ask the harder question: which next marks are allowed?

\section*{Coda}
\markboth{CODA}{CODA}

It is worth pausing over how little has been granted.

There is no road in the device yet. No wheel. No axle. No tire with a known
circumference. No second hand ticking beside the car. No coordinate system in
which the car can be placed. No real number waiting to receive a reading. No
unit called a mile, no hour over which the mile can be divided, no law telling
a count how to become a speed. There is no radar beam, no satellite clock, no
field equation, no probability measure, no geometry, no particle, and no
physics.

The room is almost empty.

That emptiness is not a failure of imagination. It is the discipline of the
chapter. Most explanations of measurement begin after a great deal has already
been supplied. They begin with a ruler, a clock, a coordinate system, a scale,
or a number line. They begin with the furniture of measurement already arranged
around the phenomenon. Chapter 1 began earlier than that. It asked what the
compiler could accept before any of that furniture had been brought in.

The answer was not dramatic. A claim could be carried with the evidence that
the claim can be decided. From such carried claims, a count could be made. The
count did not float free of its receipts. Each new mark arrived with what made
it a mark. Comparison then became possible, but comparison was not merely a
matter of length. The facts inside the count could change how direction was
read. Finally, a process could hold a mark, hold a count, and apply a rule for
what comes next.

None of this required the compiler to understand what it was for.

That is the quiet strangeness of the chapter. Lean was not asked to measure a
car. It was not asked to understand motion. It was not asked to know that a
toothed ring might turn past a Hall-effect sensor, or that a pulse might be
counted, or that a dashboard might later show forty-two miles per hour. It was
asked smaller questions. Does this definition fit the context already accepted?
Can this field be supplied? Does this term have the type claimed for it? Can
the present declaration close without borrowing a promise from the future?

The compiler answered those questions locally. It did not see the destination.
It did not need to. Each accepted object made the next question available. A
loose claim was not enough, so the device asked for a claim with its receipt.
Once such a mark could be carried, the device could ask whether marks could be
counted. Once counts could be carried, it could ask how counts compare. Once
comparison had direction, it could ask what sort of process carries a mark
forward.

This is why the wheel-speed sensor has been useful all along. The sensor does
not know speed either. It does not know the driver's intention, the weight of
the car, the state of the road, the destination, or the meaning of the number
that will eventually appear on the dashboard. Its world is narrower. A tooth
passes, or it does not. A pulse is registered, or it is not. The current count
is held. A rule says how the next accepted event changes the count.

That rule is not speed. The pulse is not speed. The count is not speed. The
displayed number on the dashboard is already the result of a later conversion,
a calibration, and a model of what wheel rotation has to do with the motion of
the car. But the pulse is the sort of thing from which speed can later be made
honest. It is small enough to be recorded and precise enough to be carried. It
does not pretend to be the phenomenon. It enters the ledger as a record.

The coda to Chapter 1 is therefore not that the compiler has learned
measurement. That would be too much. The compiler has accepted a shape that
measurement will need. It has accepted the possibility of a witnessed mark, a
record made from such marks, a direction in which records can be compared, and
a process that holds a present value while waiting for what comes next. This is
not yet an instrument in the physical sense. But it is no longer mere syntax
either. It is the first formal outline of something that can remember.

Memory is doing quiet work here. A mark that cannot be carried vanishes as soon
as it is made. A count that cannot remember its previous marks is only a
gesture. A process that cannot hold a value between events cannot behave like
an instrument. The chapter has not solved the world, but it has found a way for
one accepted distinction to survive long enough for another question to be
asked.

That survival is the beginning of the ledger.

By now the ledger is pressure, not vocabulary. Once a pulse has been counted,
the next count is not free to pretend it did not happen. Once a claim has
entered with its receipt, later claims inherit the responsibility of being
checked in the world that acceptance created. This is where the compiler and
the instrument begin to resemble each other. Neither needs to know the whole
world. Each needs a rule for accepting the next mark without losing the marks
already accepted.

There is one more irony to hold lightly. These objects are abstract enough to
fit inside a proof assistant, but they are not floating outside matter. The
definitions were checked by a compiler running on a physical machine. The
accepted terms occupied memory. The first mark lived, for a while, in silicon.
Even the most austere formal receipt had to be carried by something.

The book will return to that fact later, when it has earned more language for
it. For now, it is enough to notice the fold in the situation. A physical
machine has accepted an abstract description of how a mark can be carried. The
description does not yet contain physics. The act of checking it was physical
all the same.

Chapter 1 does not end with a law of nature. It ends with a carried mark. That
is less grand, and more useful. A law can only be tested by something capable
of keeping a record. Before the book can ask what the world does, it has to ask
what a record is allowed to remember.

The first answer is small. A mark can be carried. A count can remember that it
has been carried. A rule can say what comes next.

The next question is not what the world is. The next question is what may be
added to the ledger.
